% ============================================================
% Notes: Integral Fourier Transform + Delta + Derivatives + Convolution
% ============================================================

\section{Integral Fourier transforms}

\subsection{Fourier transform pair (physics convention)}
Let $f(x)$ be integrable (or in a suitable distribution sense). Define the Fourier transform:
\begin{equation}
\tilde f(k) = \mathcal{F}\{f\}(k)
:= \int_{-\infty}^{\infty} f(x)\,e^{-ikx}\,dx.
\end{equation}
The inverse transform is
\begin{equation}
f(x) = \mathcal{F}^{-1}\{\tilde f\}(x)
:= \frac{1}{2\pi}\int_{-\infty}^{\infty} \tilde f(k)\,e^{ikx}\,dk.
\end{equation}

\paragraph{Interpretation.}
The Fourier transform expands $f(x)$ in the continuous basis of plane waves $e^{ikx}$.

% ------------------------------------------------------------

\subsection{Basic transform properties}
\paragraph{Linearity.}
\begin{equation}
\mathcal{F}\{af+bg\} = a\tilde f + b\tilde g.
\end{equation}

\paragraph{Shift in $x$ (translation).}
If $g(x)=f(x-x_0)$ then
\begin{equation}
\tilde g(k) = e^{-ikx_0}\tilde f(k).
\end{equation}

\paragraph{Modulation (shift in $k$).}
If $g(x)=e^{ik_0 x}f(x)$ then
\begin{equation}
\tilde g(k) = \tilde f(k-k_0).
\end{equation}

\paragraph{Scaling.}
If $g(x)=f(ax)$ with $a\neq 0$ then
\begin{equation}
\tilde g(k)=\frac{1}{|a|}\tilde f\!\left(\frac{k}{a}\right).
\end{equation}

% ------------------------------------------------------------

\section{Dirac delta function and its Fourier representation}

\subsection{Defining property}
The Dirac delta $\delta(x)$ is defined by
\begin{equation}
\int_{-\infty}^{\infty}\delta(x-a)\,f(x)\,dx = f(a),
\end{equation}
for suitable test functions $f$.

\subsection{Fourier representation of $\delta$}
Start from the inverse transform of $\tilde f(k)=1$:
\begin{equation}
\delta(x) = \frac{1}{2\pi}\int_{-\infty}^{\infty} e^{ikx}\,dk.
\end{equation}
More generally,
\begin{equation}
\delta(x-x_0) = \frac{1}{2\pi}\int_{-\infty}^{\infty} e^{ik(x-x_0)}\,dk.
\end{equation}

\subsection{Delta as completeness / orthogonality of plane waves}
Using the transform pair,
\begin{equation}
\int_{-\infty}^{\infty} e^{ikx}\,dx = 2\pi\,\delta(k),
\end{equation}
and equivalently,
\begin{equation}
\int_{-\infty}^{\infty} e^{i(k-k')x}\,dx = 2\pi\,\delta(k-k').
\end{equation}

\subsection{Fourier transform of a delta}
\begin{equation}
\mathcal{F}\{\delta(x-x_0)\}(k)
= \int_{-\infty}^{\infty}\delta(x-x_0)e^{-ikx}\,dx
= e^{-ikx_0}.
\end{equation}

% ------------------------------------------------------------

\section{Fourier transforms of derivatives}

Assume boundary terms vanish (e.g.\ $f(x)\to 0$ sufficiently fast as $|x|\to\infty$),
or interpret in the distribution sense.

\subsection{First derivative}
\begin{align}
\mathcal{F}\{f'(x)\}(k)
&= \int_{-\infty}^{\infty} f'(x)e^{-ikx}\,dx\\
&= \Big[f(x)e^{-ikx}\Big]_{-\infty}^{\infty}
+ik\int_{-\infty}^{\infty}f(x)e^{-ikx}\,dx\\
&= ik\,\tilde f(k).
\end{align}
So:
\begin{equation}
\boxed{\mathcal{F}\{f'\}(k)=ik\,\tilde f(k)}.
\end{equation}

\subsection{Second derivative}
\begin{equation}
\mathcal{F}\{f''\}(k)=(ik)^2\tilde f(k)=-k^2\tilde f(k).
\end{equation}

\subsection{$n$th derivative}
\begin{equation}
\boxed{\mathcal{F}\left\{\frac{d^n f}{dx^n}\right\}(k)=(ik)^n\tilde f(k)}.
\end{equation}

\paragraph{Key use in PDEs.}
Differential operators become algebraic multipliers in Fourier space:
\[
\frac{d}{dx}\to ik,\qquad \frac{d^2}{dx^2}\to -k^2.
\]

% ------------------------------------------------------------

\section{Convolution theorem}

\subsection{Definition of convolution}
For functions $f,g$ define their convolution
\begin{equation}
(f*g)(x) := \int_{-\infty}^{\infty} f(x-t)\,g(t)\,dt
= \int_{-\infty}^{\infty} f(t)\,g(x-t)\,dt.
\end{equation}

\subsection{Fourier transform of a convolution}
\begin{equation}
\boxed{\mathcal{F}\{(f*g)(x)\}(k)=\tilde f(k)\,\tilde g(k)}.
\end{equation}

\subsection{Proof sketch (one-line computation)}
\begin{align}
\mathcal{F}\{f*g\}(k)
&= \int dx\,e^{-ikx}\int dt\,f(x-t)g(t)\\
&= \int dt\,g(t)\int du\,f(u)e^{-ik(u+t)}\\
&= \left(\int du\,f(u)e^{-iku}\right)\left(\int dt\,g(t)e^{-ikt}\right)\\
&= \tilde f(k)\tilde g(k).
\end{align}

\subsection{Inverse form: product becomes convolution}
If $h(x)=f(x)g(x)$ then
\begin{equation}
\boxed{\tilde h(k)=\mathcal{F}\{fg\}(k)=\frac{1}{2\pi}(\tilde f * \tilde g)(k)}.
\end{equation}

\subsection{Delta as identity under convolution}
Using $\delta$,
\begin{equation}
(f*\delta)(x)=\int f(x-t)\delta(t)\,dt=f(x).
\end{equation}

% ------------------------------------------------------------

\section{High-yield summary}

\begin{itemize}
\item Transform pair:
\[
\tilde f(k)=\int_{-\infty}^{\infty}f(x)e^{-ikx}dx,\qquad
f(x)=\frac{1}{2\pi}\int_{-\infty}^{\infty}\tilde f(k)e^{ikx}dk.
\]
\item Delta representation:
\[
\delta(x)=\frac{1}{2\pi}\int_{-\infty}^{\infty}e^{ikx}dk,
\qquad
\int e^{i(k-k')x}dx=2\pi\delta(k-k').
\]
\item Derivatives:
\[
\mathcal{F}\{f^{(n)}\}(k)=(ik)^n\tilde f(k).
\]
\item Convolution theorem:
\[
\mathcal{F}\{f*g\}=\tilde f\,\tilde g,\qquad
\mathcal{F}\{fg\}=\frac{1}{2\pi}(\tilde f*\tilde g).
\]
\end{itemize}
